% =====================================================================
%  Partie 3 -- Transferts de vapeur
% =====================================================================
\part{Transferts de vapeur}

\chapter{Transferts de vapeur}
\label{chap:vapeur}

\section{Introduction}

\subsection*{La vapeur d'eau}
L'eau peut passer d'un état à l'autre :
\begin{itemize}
  \item solide $\rightarrow$ liquide : fusion ; liquide $\rightarrow$ solide : solidification ;
  \item liquide $\rightarrow$ gaz : vaporisation ; gaz $\rightarrow$ liquide : liquéfaction ;
  \item solide $\rightarrow$ gaz : sublimation ; gaz $\rightarrow$ solide : condensation.
\end{itemize}
% [Figure : diagramme des changements d'état]

\subsection*{Le point triple de l'eau}
Le point triple de l'eau se situe à 0,01\,°C et environ 6\,mbar. À la pression atmosphérique (1013\,mbar), le point de fusion est à 0\,°C et le point d'ébullition à 100\,°C.
% [Figure : diagramme de phase de l'eau]

\subsection*{L'évaporation}
L'eau s'évapore au contact de l'air sans être à 100\,°C.
On estime en extérieur l'évaporation d'un plan d'eau entre 1\,mm et 10\,mm par jour en fonction des conditions météorologiques\footnote{M.E. Jensen, \emph{Estimating evaporation from water surfaces}, 2010.}.

Pour s'évaporer, l'eau doit changer d'état.
Elle absorbe l'énergie thermique de sa surface hôte pour s'évaporer.
Cette énergie est la chaleur latente.

\subsection*{La condensation}
% [Figures : bouteille froide, rosée, vitre, moisissures]

\section{Pression de vapeur et humidité}

\subsection*{Pression de vapeur}
Dans le cas de l'air humide, on appelle pression de vapeur, $p_v$, la pression (dite partielle) de la vapeur contenue dans l'air.
La pression de vapeur reflète la proportion de vapeur d'eau contenue dans l'air.

\subsection*{Pression de vapeur saturante}
La pression de vapeur saturante, $p_{vs}$, est la pression à laquelle la phase gazeuse est en équilibre avec sa phase liquide ou solide à une température donnée.
À l'équilibre, l'eau se rassemblera sous forme de gouttelettes sphériques pour équilibrer les forces en présence.
\begin{itemize}
  \item C'est la frontière à laquelle un gaz se changera en liquide ou en solide.
  \item La pression de vapeur saturante dépend de la température.
\end{itemize}

\subsection*{Humidité relative}
L'humidité relative de l'air, HR, ou degré hygrométrique, correspond au rapport de la pression partielle de vapeur d'eau contenue dans l'air sur la pression de vapeur saturante à la même température.

L'humidité relative est analogue à une mesure du rapport entre le contenu en vapeur d'eau de l'air et sa capacité maximale à en contenir dans ces conditions de température.
Ce rapport changera si on change la température.
L'humidité relative s'exprime en pourcentage :
\begin{equation}
  HR = \frac{p_v}{p_{vs}(T)} \times 100
\end{equation}
$p_{vs}$ peut se trouver dans des abaques dans lesquels elle est déjà calculée ou mesurée pour plusieurs températures.

Plus l'air est chaud, plus il peut contenir d'eau sous forme de vapeur.
Inversement, quand l'air se refroidit, la vapeur se condense et forme des gouttelettes d'eau liquide : on dit que le seuil de saturation augmente avec la température.

\subsection*{Diagramme de l'air humide}
Le diagramme de l'air humide permet de relier la température à la pression de vapeur saturante.
Par exemple, l'air à 20\,°C et 50\,\% d'humidité relative condense autour de 9\,°C.
% [Figures : diagramme de l'air humide (pression atmosphérique 101325 Pa), lecture du point de rosée, zone de confort]

\section{Passage de la vapeur}

\subsection*{Matériaux perspirants}
Les matériaux perspirants sont des matériaux qui permettent la diffusion de la vapeur d'eau.

Chaque matériau offre une certaine résistance à la migration de la vapeur d'eau ; cette résistance s'exprime par :
\begin{itemize}
  \item le coefficient de résistance à la diffusion de vapeur d'eau : $\mu$ (sans unité) ;
  \item la résistance à la diffusion de vapeur d'eau : $S_d$ (en mètres).
\end{itemize}

\subsection*{Coefficient de résistance à la diffusion de vapeur d'eau}
Le coefficient de résistance à la diffusion de vapeur $\mu$ représente la capacité d'un matériau à s'opposer au passage de la vapeur.

C'est une grandeur relative résultant du rapport entre la perméabilité à la vapeur d'eau de l'air et celle du matériau.
Plus $\mu$ est élevé, plus la résistance est élevée ($\mu$ est infini pour un métal et vaut 1 pour l'air).

La détermination du facteur de résistance à la diffusion de la vapeur dans un matériau est réalisée suivant un protocole défini dans la norme NF EN ISO 12572.

\subsection*{Résistance à la diffusion de vapeur d'eau}
La quantité de vapeur d'eau diffusant à travers une couche de matériau dépend aussi de l'épaisseur.
On appelle alors épaisseur équivalente de diffusion, $S_d$, la résistance qu'offre une couche de matériau à la diffusion de vapeur pour une épaisseur donnée.

$S_d$ est le produit de l'épaisseur $e$ d'un matériau et du coefficient de résistance à la diffusion de vapeur $\mu$ :
\begin{equation}
  S_d\ \mathrm{[m]} = \mu \times e\ \mathrm{[m]}
\end{equation}

Le $S_d$ d'une couche de matériau correspond à l'épaisseur en mètres de la couche d'air stationnaire qui exercerait la même résistance à la diffusion de vapeur que la couche de matériau.
Exemple : lorsque le $S_d$ d'un matériau vaut 5, cela signifie que cette couche de matériau équivaut à 5 mètres d'air stationnaire.

Le coefficient de résistance à la diffusion de vapeur de certains matériaux dépend de leur état : sec ou humide.

\subsection*{Frein-vapeur, pare-vapeur et pare-pluie}
Le \textbf{frein-vapeur} est un film fin qui freine le passage de la vapeur d'eau (murs perspirants). Il aura un $S_d$ de 1 à 5\,m.

On parle de frein-vapeur « intelligent » lorsque la membrane du frein-vapeur est à diffusion variable. Elle est alors composée de pores qui s'ouvrent ou se ferment en fonction de l'humidité de l'air : elle peut être frein-vapeur ou pare-vapeur.

Le \textbf{pare-vapeur} empêche le passage de la vapeur d'eau ; il se place du côté intérieur de l'habitation et possède un $S_d > 10$\,m.

Remarque : le \textbf{pare-pluie} se posera sous toiture ou à l'extérieur du mur. Son rôle est d'assurer l'étanchéité à l'eau, mais il doit laisser passer la vapeur d'eau.
Le $S_d$ d'un pare-pluie est de l'ordre de 0,18\,m.

\subsection*{Condensation de surface}
L'air intérieur peut arriver en contact avec une surface de température plus faible. Au contact, l'air se refroidit et la température qu'il atteint dans le voisinage de la surface peut se retrouver en dessous du point de rosée.
Il apparaît alors de la condensation dite \textbf{de surface} (la pression de vapeur dans l'air atteint la pression de saturation).
% [Figures : moisissures en angle]

Pour empêcher la condensation de surface, il faut :
\begin{itemize}
  \item avoir une qualité suffisante de l'isolation de l'enveloppe, de façon à ce que la température de la face intérieure de celle-ci ne descende pas en dessous de 17,5\,°C ;
  \item diminuer, par une ventilation complémentaire, l'humidité relative de l'air intérieur ;
  \item chauffer la face intérieure des parois froides.
\end{itemize}
Si ces mesures ne sont pas applicables ou sont insuffisantes, la seule façon de limiter les dégâts est de traiter la surface des parois pour empêcher la pénétration de l'eau de condensation par capillarité dans les parois, et de pouvoir recueillir et évacuer l'eau condensée.

\subsection*{Condensation interne}
La condensation interne se produit si, à un endroit d'un élément, la pression de vapeur réelle devient égale à la pression de saturation régnant à cet endroit.
Ce phénomène résulte des différences de pression de vapeur et de température.

Il existe bien souvent une différence de pression de vapeur entre l'intérieur et l'extérieur (production de vapeur par les occupants, les plantes, la lessive, la cuisson, le nettoyage).

En hiver, la pression partielle de vapeur intérieure est, en général, supérieure à celle correspondant au climat extérieur, et on a alors un flux de vapeur qui se produit.

Fonctionnement :
\begin{itemize}
  \item En hiver, la pression de vapeur est souvent plus élevée à l'intérieur qu'à l'extérieur.
  \item Ce différentiel de pression de vapeur engendre une migration de vapeur par diffusion vers l'extérieur.
  \item Au fur et à mesure qu'elle traverse les différents matériaux constituant l'enveloppe, la vapeur se rapproche de l'extérieur et se refroidit progressivement.
  \item Si la température du point de rosée est atteinte, la vapeur se condense : on parlera de condensation d'hiver.
  \item L'humidité peut dégrader les matériaux et avoir des conséquences sur la durabilité de la paroi et de ses performances, ainsi que sur le confort et la santé des habitants.
\end{itemize}
Contrairement à la condensation de surface, la condensation interne n'est pas visible directement.

Pour empêcher la condensation interne :
\begin{itemize}
  \item si on veut éviter la condensation interne dans une paroi constituée de plusieurs couches de matériaux différents, il faut que la perméabilité à la vapeur de ceux-ci augmente de l'intérieur vers l'extérieur ;
  \item si ce n'est pas possible (par exemple avec certaines techniques d'isolation par l'intérieur), il faut poser un pare-vapeur du côté chaud de la paroi pour provoquer une chute de la pression de vapeur avant l'isolant.
\end{itemize}

\subsection*{Condensation d'été}
\begin{itemize}
  \item En été et au printemps, la température et l'humidité relative de l'air extérieur sont parfois plus élevées qu'à l'intérieur : la pression de vapeur peut donc être plus élevée à l'extérieur qu'à l'intérieur.
  \item L'humidité présente dans les matériaux de la paroi a alors tendance à migrer vers l'intérieur.
  \item Si elle est bloquée par une éventuelle membrane, la vapeur qui migre vers l'intérieur ne peut alors plus s'évaporer.
  \item Au contraire, elle peut condenser en arrivant contre celle-ci (au final, c'est le potentiel de séchage du mur qui est affaibli, et l'humidité risque alors de s'accumuler).
\end{itemize}
On parle alors de condensation d'été. Elle apparaît entre la membrane et l'isolant, provoquant une perte de performance de l'isolant humidifié et des risques de moisissures s'il y a du bois ou des matériaux organiques dans cette couche du mur.

Pour éviter ce problème, on place des membranes dites « intelligentes », aussi appelées à $\mu$ variable, qui permettent de réduire ce type de risque.

\section{Évolution de la pression de vapeur dans une paroi}

\subsection*{Paroi simple}
Pour une paroi simple, considérons la diffusion de vapeur à travers une couche de matériau plane et homogène d'épaisseur $e$.
Cette diffusion résulte d'une différence de pression partielle de vapeur.
On note :
\begin{itemize}
  \item $A$ la superficie (m$^2$) ;
  \item $p_{v1}$ la pression de vapeur à la face intérieure (Pa) ;
  \item $p_{v2}$ la pression de vapeur à la face extérieure (Pa) ;
  \item $p_{v1} > p_{v2}$ (conditions hivernales).
\end{itemize}
% [Figure : profil de pression de vapeur dans une paroi simple]

La différence de pression de vapeur $(p_{v1} - p_{v2})$ entraîne alors une diffusion de la vapeur d'eau de l'intérieur vers l'extérieur.
La densité de flux de vapeur traversant la paroi sera :
\begin{equation}
  q_d = \frac{p_{v1} - p_{v2}}{5{,}4\times10^{9}\,\mu\,e} \qquad \mathrm{[kg/m^2.s]}
\end{equation}
$5{,}4\times10^{9}\,\mu\,e$ est appelé « la résistance à la diffusion » et est noté $Z$.
On peut donc réécrire l'équation de $q_d$ sous la forme :
\begin{equation}
  q_d = \frac{p_{v1} - p_{v2}}{Z}
\end{equation}
On pourra noter le parallèle avec la diffusion de la chaleur dans une paroi :
\begin{equation}
  \varphi = U\,\Delta T = \frac{T_1 - T_2}{R}
\end{equation}

On peut donc obtenir une résistance à la diffusion très élevée en utilisant soit une couche mince d'un matériau ayant une valeur $\mu$ très élevée (= pare-vapeur), soit une couche épaisse d'un matériau ayant une valeur $\mu$ peu élevée.

\paragraph{Exemple}
Une paroi en briques silico-calcaires (1800\,kg/m$^3$) a une valeur $\mu$ de 15 et une épaisseur de 20\,cm. Si la pression de vapeur à l'extérieur est de 280\,Pa et celle à l'intérieur de 1320\,Pa, on aura alors la densité de flux de vapeur :
\begin{equation}
  q_d = \frac{p_{v1} - p_{v2}}{5{,}4\times10^{9}\,\mu\,e}
      = \frac{1320 - 280}{5{,}4\times10^{9}\times 15\times 0{,}2}
\end{equation}
\begin{equation}
  q_d = 64{,}2\times10^{-9}\,\mathrm{kg/m^2.s} = 5{,}5\,\mathrm{g/m^2.jour}
\end{equation}
% NB : la diapo indiquait 15 cm d'épaisseur mais le calcul utilise 0,2 m.
% Avec e = 0,15 m : q_d = 85,6e-9 kg/m2.s = 7,4 g/m2.jour.

\subsection*{Paroi composée}
Pour une paroi composite, la résistance à la diffusion totale s'obtiendra en faisant la somme des résistances à la diffusion de chacune des couches :
\begin{equation}
  Z_t = Z_1 + Z_2 + \dots + Z_n
\end{equation}
ou encore :
\begin{equation}
  Z_t = 5{,}4\times10^{9}\left(\mu_1 e_1 + \mu_2 e_2 + \dots + \mu_n e_n\right)
\end{equation}
La formule de la densité de flux de vapeur devient donc tout simplement :
\begin{equation}
  q_d = \frac{p_{v1} - p_{v2}}{5{,}4\times10^{9}\,\sum \mu\,e} \qquad \mathrm{[kg/m^2.s]}
\end{equation}
% [Figure : paroi composée -- enduit extérieur, isolation, maçonnerie, enduit intérieur]

En régime stationnaire, les quantités $q_{d1}$, $q_{d2}$, $q_{d3}$ et $q_d$ doivent être égales.
En effet, si $q_{d1}$ était supérieur à $q_{d2}$, la quantité de vapeur d'eau quittant le plan 1 serait supérieure à celle qui y arrive et, dans ce cas, la pression de vapeur au droit du plan 1 ($p_{v1}$) ne pourrait rester constante.

Tout plan parallèle est donc traversé par la même quantité d'humidité : $q_d = q_{d1} = q_{d2} = q_{d3}$.
On peut voir ici un parallèle évident avec les transferts de chaleur.

Comme pour les transferts de chaleur, on peut appliquer la formule de la densité de flux de vapeur (pour 3 matériaux, par exemple) :
\begin{equation}
  q_d = \frac{p_{vi} - p_{ve}}{5{,}4\times10^{9}\left(\mu_1 e_1 + \mu_2 e_2 + \mu_3 e_3\right)}
\end{equation}
On a donc :
\begin{equation}
  q_{d1} = \frac{p_{vi} - p_{v12}}{5{,}4\times10^{9}\,\mu_1 e_1} \qquad
  q_{d2} = \frac{p_{v12} - p_{v23}}{5{,}4\times10^{9}\,\mu_2 e_2} \qquad
  q_{d3} = \frac{p_{v23} - p_{ve}}{5{,}4\times10^{9}\,\mu_3 e_3}
\end{equation}
On peut donc écrire que :
\begin{equation}
  q_d = \frac{p_{vi} - p_{ve}}{5{,}4\times10^{9}\,\sum \mu\,e}
      = \frac{p_{vi} - p_{v12}}{5{,}4\times10^{9}\,\mu_1 e_1}
      = \frac{p_{v12} - p_{v23}}{5{,}4\times10^{9}\,\mu_2 e_2}
      = \frac{p_{v23} - p_{ve}}{5{,}4\times10^{9}\,\mu_3 e_3}
\end{equation}
D'où l'on peut extraire facilement les pressions aux interfaces :
\begin{equation}
  p_{v12} = p_{vi} - \frac{p_{vi} - p_{ve}}{\sum \mu\,e}\,\mu_1 e_1
\end{equation}
\begin{equation}
  p_{v23} = p_{v12} - \frac{p_{vi} - p_{ve}}{\sum \mu\,e}\,\mu_2 e_2
\end{equation}

On pourra donc, en fonction des besoins, tracer différents profils de température et d'humidité permettant de vérifier s'il y a, ou non, condensation dans les murs.

Tant que $p_v$ est inférieure à $p_{vs}$, tout va bien : il n'y a pas de risque de condensation dans le mur. En revanche, dès qu'on la dépasse, il y a condensation.

On peut tracer l'évolution de $p_v$, mais également du rapport de $p_v$ sur $p_{vs}$ pour tracer l'évolution de l'humidité relative (HR) au sein de la paroi.
C'est l'évolution de HR qui est représentée sur Ubakus\footnote{Ubakus : \texttt{https://www.ubakus.de/fr/calculateur-valeur-u/}}.

Étant donné que la pression de vapeur saturante dépend de la température, il faut au préalable calculer l'évolution de la température au sein de la paroi.

Lorsque l'on conçoit une typologie de paroi, il faut vérifier que la condensation ne se produit pas à l'intérieur de celle-ci dans différentes conditions (hiver et été).
